% Written last and revised after peer review. The earlier version quoted a
% regression trend whose p-value assumed an independence that overlapping
% five-year windows rule out, and described the amplitude comparison as an
% exact reproduction when ten observations cannot establish that. Both are
% restated here.
\noindent
Every large dealer reporting market risk under Item 7A publishes Value-at-Risk
by risk category and a single reconciling line between those categories and the
firmwide total. The category figures are levels. The reconciling line, the diversification
effect, is a relationship, and the relationship it reports is the one thing
the filing withholds. That asymmetry makes it the single disclosed quantity an
outside party can check. This paper tests it on ten consecutive Goldman Sachs filings,
fiscal years 2016 through 2025, with a historical-simulation engine run at the
bank's own disclosed settings --- full revaluation, a five-year scenario window
and age-weighted scenarios --- over a proxy book of seventeen public risk
factors sized to the disclosed split of risk across categories.

The level reproduces and the steadiness does not. Goldman's effect has averaged
42.0\% of summed category VaR with a year-to-year standard deviation of 2.08
percentage points; a book sized once and then left alone averages 41.6\%, which
no test on ten annual observations separates from the disclosure, with a
standard deviation of 6.33 ($F = 9.26$, $p = 0.003$). The firm's second
disclosed measure locates the failure. Period-end against period-end the two
amplitudes cannot be told apart; year on year the disclosed series carries
nothing forward (lag-1 autocorrelation $-0.47$) while the proxy carries almost
everything ($+0.9$). What fails to reproduce is not how far the ratio moves but
whether it comes back. Sweeping the undisclosed decay rate across
sixteen values, from equal weighting to a 6.6-day half-life, never brings the
proxy nearer than three times the disclosed steadiness, and the rates that best
repair the engine's coverage are among the worst on steadiness. Across all
sixteen sleeve-sign configurations, five match the disclosed steadiness and none
of those averages above 35.6\%, eight reach 40\% and the steadiest of those
varies by 3.68 points; none does both.
